\section{Flattening the octree neighbourhood}
\label{sec:flatten}

\subsection{What the shader needs}

A chunk's scan cell can query a neighbouring cell at a coarser or finer
LoD, possibly in a neighbouring chunk.  The shader cannot walk pointers, so
the host must provide a node table that contains every cell the walk can
reach.  The table must include:

\begin{enumerate}
\item the world root at index 0 (the shader's descent anchor);
\item the root-to-chunk ancestor path, with root-consistent cubes;
\item all descendants of the chunk down to the target LoD;
\item all coarser nodes in the one-cell margin around the chunk, plus
  their children, because an edge sample can cross the boundary and land
  in a coarser neighbour;
\item a child index per node (or $-1$ when absent);
\item one dense vertex slot per surface node in the table.
\end{enumerate}

\begin{figure}[htbp]
\centering
\begin{tikzpicture}[scale=0.9]
  % chunk grid
  \draw[grid] (0,0) rectangle (5,5);
  \draw[grid] (2,0)--(2,5); \draw[grid] (3.5,0)--(3.5,5);
  \draw[grid] (0,2)--(5,2); \draw[grid] (0,3.5)--(5,3.5);
  \fill[accent!12] (2,2) rectangle (3.5,3.5);
  \node[font=\scriptsize] at (2.75,2.75) {chunk};
  % margin
  \draw[accent2,thick,dashed] (-0.7,-0.7) rectangle (5.7,5.7);
  \node[font=\scriptsize,accent2] at (5.2,5.55) {one-cell margin};
  % root path
  \draw[good,thick] (0,0)--(0.8,0.8)--(2,2);
  \node[font=\scriptsize,good] at (0.7,0.25) {root path};
  % neighbor coarse
  \fill[orange!15] (3.5,3.5) rectangle (5,5);
  \node[font=\scriptsize] at (4.25,4.25) {coarse};
  % arrows
  \draw[-{Latex[length=2mm]},ink!60] (2.75,1.3) -- (2.75,2.0);
  \node[right,font=\scriptsize] at (2.8,1.55) {scan cell};
\end{tikzpicture}
\caption{The flattened neighbourhood.  The chunk subtree is walked to the
target rung; the one-cell margin is included so edge samples that cross
the boundary can resolve coarser neighbours.  Index 0 is always the world
root, because the shader's descent starts there for samples outside the
chunk anchor cube.}
\label{fig:flatten}
\end{figure}

\subsection{Root-consistent ancestor chain}

The host replays the child index from the world root to the chunk:
\begin{lstlisting}[caption={Ancestor chain replay.},label={lst:chain}]
cube = rootCube; node = root;
chain.push({node, cube});
while (node != chunk) {
    i = cube.getChildIndex(chunkCenter);
    node = children[i];
    cube = cube.getChild(i);
    chain.push({node, cube});
}
\end{lstlisting}
The replay is necessary because a cached chunk cube may have accumulated
float drift.  The shader will compute its own child indices from the world
root, so the host must use exactly the same arithmetic.  If the chain
cannot be replayed (stale node, missing child), the chunk is marked
\code{Unavailable}; there is no CPU fallback.

\subsection{Descending the table}

The table is built from the chain root.  For each node, the host writes:
\begin{center}
\begin{tabular}{@{}ll@{}}
\toprule
Buffer & Contents\\
\midrule
\code{sdf0, sdf1} & two \code{vec4}s of corner samples\\
\code{cubeMinSize} & \code{vec4(cubeMin, cubeSize)}\\
\code{meta} & \code{ivec4(lod, type, brush, flags)}\\
\code{hsv} & \code{vec4(hsv, textureType)}, \code{vec4(data,0,0)}\\
\code{child03, child47} & two \code{ivec4}s of child indices\\
\code{cells} & compacted scan-cell node ids\\
\bottomrule
\end{tabular}
\end{center}
The layout is 128 bytes per node: $4\times4$ floats for SDFs, $4$ floats
for the cube, $4$ ints for metadata, $8$ floats for HSV/mapping, and $8$
ints for children.  It is uploaded into host-visible coherent buffers; the
shader reads it through storage-buffer descriptors.

\subsection{The scan-cell list}

The scan list contains cells exactly at the target rung that have at least
one crossing edge.  The host counts the crossing edges per cell while it
builds the list:
\begin{equation}
  \text{crossEdges}(c) =
  \#\{(i,j)\in\text{SDF\_EDGES} :
      \operatorname{sign}(d_i)\ne\operatorname{sign}(d_j)\},
\end{equation}
with $+\infty$ counting as non-negative.  This is not the same as
\code{type == Surface}: a negative-plus-infinity cell is typed Solid yet
can emit an edge.  Filtering by type would silently drop those edges.

The same pass produces the index-capacity hint used by the geometry
dispatch:
\begin{equation}
  I_{\text{hint}} =
  6 \sum_{c\in\text{cells}}
      \operatorname{clamp}\!\left(
        \frac{\text{cellSize}(c)}{m},1,64\right)
      \cdot\text{crossEdges}(c),
\end{equation}
where $m$ is the smallest cube in the flattened table.  A segment expands
to at most two triangles of three indices; the clamp is the shader's break
list capacity.  This is a topology bound, not a statistical estimate.  If
it is nevertheless exceeded, the shader reports the exact requirement and
the host grows only that chunk's span.

\subsection{The descent anchor}

The shader's \code{findCellAt} descends from the root.  A per-sample root
descent would be expensive, so the host supplies \code{fromId}: the first
node on the chunk path whose stored LoD is at most the target LoD.  For any
sample inside that node's cube, starting the descent there is
bit-identical to starting at the root, because no ancestor of the anchor
is a stop node.  Samples in the margin outside the anchor fall back to the
root.  This is a pure optimization; the result is the same.

\subsection{Why the margin is one cell}

An edge sample point is a corner of a scan cell, nudged outward by one ulp
in two axes.  The neighbouring quadrant cell can be one cell larger than
the scan cell (a coarse ancestor).  The flattening domain is therefore
\begin{equation}
  [\mathbf{c}_{\min}-\ell,\; \mathbf{c}_{\max}+\ell],
\end{equation}
where $\ell$ is the largest target-rung cell size in the subtree,
\emph{including non-surface cells that were compacted out of the scan
list}.  Deriving the margin from the compacted list would under-estimate
the largest neighbour and could omit a cell the walk samples.

\subsection{Vertex slots}

Every node whose \code{type} is \code{Surface} receives a dense slot in
table order:
\begin{equation}
  \text{flags} \mathrel{|}= (s+1)\ll1.
\end{equation}
Bit 0 remains the leaf flag.  The shader decodes
\begin{equation}
  s = (\text{flags}\gg1)-1.
\end{equation}
The count is the dense slot address space, an exact upper bound on every
index the shader can write.  It is not the number of vertices actually
produced: the geometry pass writes the payload only for slots referenced by
a triangle, and the published span is trimmed to the highest referenced
slot plus one.  Unreferenced slots between zero and that extent remain holes
in the reserved span, so the per-chunk profile reports both the written
vertex count and the published extent.

\subsection{Table order and determinism}

The table order is fixed by the traversal: root, chunk subtree in
child-index order, then the margin descent.  The slot assignment follows
that order.  A rebuild with the same tree produces the same slots.  The
renderer relies on this for the packed-pool span and for the ray-tracing
BLAS, which references the same base vertex and index ranges.

\subsection{Failure modes}

Flattening can fail for only a few reasons:
\begin{itemize}
\item the chunk node is not reachable from the root (stale parameters);
\item the chain is deeper than 200 (a corrupt tree);
\item the chunk has no crossing scan cell, or no surface node in the
  flattened table owns a slot (a valid empty chunk, not a failure);
\item the table is missing a child block the descent needs.
\end{itemize}
The first, second, and fourth mark the chunk \code{Unavailable}; the
caller skips it and logs.  The third returns a valid empty input, the
geometry pass is not dispatched, and the publisher settles the empty
revision.
